\documentclass[11pt,twoside]{amsart}
\usepackage{amssymb, amsmath, enumerate, palatino, hyperref}
\usepackage[normalem]{ulem}
\usepackage{fullpage}
\usepackage[T1]{fontenc}
\renewcommand{\labelitemi}{\guillemotright}
\usepackage{mathrsfs}

\usepackage{tikz}
\tikzstyle{ball} = [circle,shading=ball, ball color=black,
    minimum size=1mm,inner sep=1.3pt]
\usetikzlibrary{shapes}

\theoremstyle{plain}
\newtheorem{prop}{Proposition}%[section]
\newtheorem{lemma}[prop]{Lemma}
\newtheorem{thm}[prop]{Theorem}
\newtheorem{obs}[prop]{Observation}
\newtheorem{app}[prop]{Application}
\newtheorem*{MainThm}{Main Theorem}
\newtheorem{cor}[prop]{Corollary}
\newtheorem{conj}[prop]{Conjecture}
\theoremstyle{remark}
\newtheorem{rmk}[prop]{Remark}
\newtheorem{prob}{Problem}
\newtheorem{bonus}[prop]{Bonus Problem}
\newtheorem{exc}{Exercise}
\theoremstyle{definition}
\newtheorem{ex}[prop]{Example}
\theoremstyle{definition}
\newtheorem{defn}[prop]{Definition}

\newcommand{\RR}{\mathbb{R}}
\newcommand{\ZZ}{\mathbb{Z}}
\newcommand{\CC}{\mathbb{C}}
\newcommand{\NN}{\mathbb{N}}
\newcommand{\QQ}{\mathbb{Q}}
\newcommand{\PP}{\mathbb{P}}

\newcommand{\cC}{\mathscr{C}}
\newcommand{\Ob}{\operatorname{Ob}}
\newcommand{\Mor}{\operatorname{Mor}}

\newcommand{\Set}{\operatorname{Set}}
\newcommand{\FinSet}{\operatorname{FinSet}}
\newcommand{\Vect}{\operatorname{Vect}}
\newcommand{\FinVect}{\operatorname{FinVect}}
\newcommand{\Mat}{\operatorname{Mat}}
\newcommand{\Gp}{\operatorname{Gp}}
\newcommand{\FinGp}{\operatorname{FinGp}}
\newcommand{\AbGp}{\operatorname{AbGp}}
\newcommand{\Ring}{\operatorname{Ring}}
\newcommand{\CommRing}{\operatorname{CommRing}}
\newcommand{\Field}{\operatorname{Field}}
\newcommand{\Top}{\operatorname{Top}}
\newcommand{\Aut}{\operatorname{Aut}}
\newcommand{\id}{\operatorname{id}}
\DeclareRobustCommand{\stir}{\genfrac\{\}{0pt}{}}

% For solutions.
\newenvironment{sol}{
  \medskip

  \noindent\textsc{Solution:}
}
{
  \medskip

}

\newcommand{\defeq}{:=}

\title{Math 113: Discrete Structures\\ Homework 22}
%\author{} %Remove the leading % and put your name here if using this .tex file as a template for your homework.

\begin{document}
\maketitle

\noindent
\textbf{Due:} Monday, March 30 at 10pm.
\bigskip

\begin{prob}
  Let~$h_n$ be the number of hexagons in a central packing of hexagons with~$n+1$
  layers.  The first few values are pictured below:
  \medskip 

  \begin{center}
    \begin{tikzpicture} [scale=0.5,hexa/.style= {shape=regular polygon,regular polygon
      sides=6,minimum size=0.5cm, draw,inner sep=0,anchor=south,fill=blue!10,rotate=30}] 
      \node[hexa] at (0,0) {};
      \node at (0,-3.5) {$h_0=1$};
      \begin{scope}[xshift=5cm]
        \node[hexa] at (0,0) {};
	\foreach \j in {0,...,5} {
	  \node[hexa] at (\j*60:0.866) {};
	}
        \node at (0,-3.5) {$h_1=7$};
      \end{scope}
      \begin{scope}[xshift=12cm]
        \node[hexa] at (0,0) {};
	\foreach \j in {0,...,5} {
	  \node[hexa] at (\j*60:0.866) {};
	  \node[hexa] at (\j*60:2*0.866) {};
	  \node[hexa] at (\j*60+30:1.5) {};
	}
	\node at (0,-3.5) {$h_2=19$};
      \end{scope}
      \begin{scope}[xshift=20cm]
        \node[hexa] at (0,0) {};
	\foreach \j in {0,...,5} {
	  \node[hexa] at (\j*60:0.866) {};
	  \node[hexa] at (\j*60:2*0.866) {};
	  \node[hexa] at (\j*60:3*0.866) {};
	  \node[hexa] at (\j*60+30:1.5) {};
	}
	\node[hexa] at (0.866*0.5,0.866*3*0.866) {};
	\node[hexa] at (-0.866*0.5,0.866*3*0.866) {};
	\node[hexa] at (0.866*0.5,-0.866*3*0.866) {};
	\node[hexa] at (-0.866*0.5,-0.866*3*0.866) {};
	\node at (0,-3.5) {$h_3=37$};
      \end{scope}
      \begin{scope}[xshift=20cm,rotate=60]
	\node[hexa] at (0.866*0.5,0.866*3*0.866) {};
	\node[hexa] at (-0.866*0.5,0.866*3*0.866) {};
	\node[hexa] at (0.866*0.5,-0.866*3*0.866) {};
	\node[hexa] at (-0.866*0.5,-0.866*3*0.866) {};
      \end{scope}
      \begin{scope}[xshift=20cm,rotate=-60]
	\node[hexa] at (0.866*0.5,0.866*3*0.866) {};
	\node[hexa] at (-0.866*0.5,0.866*3*0.866) {};
	\node[hexa] at (0.866*0.5,-0.866*3*0.866) {};
	\node[hexa] at (-0.866*0.5,-0.866*3*0.866) {};
      \end{scope}
    \end{tikzpicture}
  \end{center}

  Recall the first difference operator~$\Delta[h]_n=h_{n+1}-h_n$.
  \begin{enumerate}[(a)]
    \item Notice that the pictures above nest in each other about the origin.
      Use that fact to draw pictures for the first differences~$\Delta[h]$ and
      use the pictures to determine the sequence of second differences~$\Delta^2[h]$. Make sure you do prove indeed your formula for the second difference.
    \item \textbf{Using the theory of difference operators from the text} and the results
      above, compute a polynomial~$p(n)$ such that~$p(n)=h_n$ for all~$n\geq0$.
  \end{enumerate}
\end{prob}

\begin{prob} Let~$F_n$ be the~$n$-th Fibonacci number.
Our text uses induction to show that
\begin{equation}\label{eqn: Fib square identity}
  F_{2n+1}=F_n^2+F_{n+1}^2.
\end{equation}
You are now asked to give a combinatorial proof using tilings of checkerboards.
Let~$a_n$ be the number of ways of tiling a~$2\times n$ checkerboard
with~$2\times 1$ dominoes.  At the end of the section on induction, our text
shows that $a_n=F_{n+1}$. 

Rewrite equation~\eqref{eqn: Fib square identity} in terms of appropriate~$a_i$
and prove the resulting (equivalent) formula by counting tilings of a~$2\times
2n$ checkerboard.  (Hint: Our checkboard has two halves, each of size~$2\times
  n$. Consider how dominoes in a tiling behave at the middle where these two
  halves meet. There are only two possibilities: there is a pair of dominoes
  straddling the two halves, or no domino straddles the center.  Given that each
  half of our chessboard contains an even number of dominoes, it is impossible
for only one domino to straddle the center.)
\end{prob}

\end{document}
